\documentclass[10pt,a4paper]{amsart}
\usepackage[margin=19mm]{geometry}
\usepackage{amsmath,amssymb,mathtools}
\usepackage[scr=rsfs,cal=euler]{mathalfa}
\usepackage{xfrac}
\usepackage[unicode,hidelinks,bookmarks=false]{hyperref}
\newcommand{\ie}{i.e.\ }
\newcommand{\cf}{cf.\ }
\newcommand{\resp}{respectively\ }
\newcommand{\Subsection}{\subsection}
\providecommand{\nobreakdash}{\nobreak\hbox{-}}
\setlength{\emergencystretch}{2em}
\multlinegap0pt
\usepackage{amssymb}
\usepackage{xcolor}
\usepackage{float}
\usepackage{tikz}
\newtheorem{prop}{Proposition}
\renewcommand\eqref[1]{(\ref{#1})} %Need with hyperref
\title{On The Fast Fourier Transform on SU(2)}
\author{Julio Delgado, Alejandro Uma\~na}
\date{Source: arXiv:2605.23923v2 (CC BY 4.0)}
\begin{document}
\maketitle

\noindent\emph{Source and licence.} On The Fast Fourier Transform on SU(2). Version arXiv:2605.23923v2 (CC BY 4.0), distributed by arXiv under the Creative Commons Attribution 4.0 International licence, http://creativecommons.org/licenses/by/4.0/, as recorded in the arXiv licence metadata for this exact version and shown on the abstract page https://arxiv.org/abs/2605.23923v2. Full licence notice: \texttt{licenses/arxiv-2605.23923-CC-BY-4.0.txt}.\par\smallskip

\noindent\textit{Editorial scope.} Counted unit: the article's prop prop (source0.tex lines 811--851). The article's complete original proof of it is delivered with it (source0.tex lines 853--943, one \texttt{proof} environment of 91 lines). No mathematics is written, repaired or re-derived: the statement and the proof are the article's own lines, in the article's own notation, and no retained line is substituted or re-flowed.  No further source-local context is needed: every cross-reference of every retained line resolves inside the two delivered spans.  Nothing was omitted from any retained span, so the package carries no omission record.  No retained line contains a citation, so the wrapper carries no bibliography and no bibliography entry.  No retained line contains a figure environment, an image-inclusion command or drawing code, so no image is delivered and the package declares no asset dependency.  Editorial formatting only, in addition: an AMS \texttt{amsart} preamble that restates the article's own declarations and macros used by the retained lines, namely the environments \texttt{prop} on separate counters, together with the standard packages those lines need.  The retained environments print in this document as Proposition 1 in order of appearance, whereas the article prints the same items with the numbering of its own full document.  The complete native source of the article as distributed in the arXiv e-print for this exact version is bundled unchanged as \texttt{sources/201247/source0.tex}.
\begin{prop}
Let $L$ be fixed. Assume that for each $r \geq 1$ the Legendre polynomials satisfy
\[
P_{L+r}(\cos\theta)
= A_r^L(\cos\theta)\, P_L(\cos\theta)
+ B_r^L(\cos\theta)\, P_{L-1}(\cos\theta),
\]
with initial conditions
\[
A_0^L = 1, \quad B_0^L = 0, \qquad
A_{-1}^L = 0, \quad B_{-1}^L = 1.
\]
Then the coefficients $A_r^L$ and $B_r^L$ satisfy the matrix recurrence
\[
\begin{pmatrix}
A_{r+1}^L & B_{r+1}^L \\
A_r^L & B_r^L
\end{pmatrix}
=
\begin{pmatrix}
\dfrac{2L+2r+1}{L+r+1}\cos\theta &
-\dfrac{L+r}{L+r+1} \\
1 & 0
\end{pmatrix}
\begin{pmatrix}
A_r^L & B_r^L \\
A_{r-1}^L & B_{r-1}^L
\end{pmatrix},
\]
or, equivalently, the scalar recurrences
\[
\begin{cases}
A_{r+1}^L
= \dfrac{2L+2r+1}{L+r+1}\cos\theta\, A_r^L
- \dfrac{L+r}{L+r+1} A_{r-1}^L, \\[0.8em]
B_{r+1}^L
= \dfrac{2L+2r+1}{L+r+1}\cos\theta\, B_r^L
- \dfrac{L+r}{L+r+1} B_{r-1}^L.
\end{cases}
\]
\end{prop}

\begin{proof}
We start from the three-term recurrence relation for Legendre polynomials,
\[
(L+r+1) P_{L+r+1}(\cos\theta)
= (2L+2r+1)\cos\theta\, P_{L+r}(\cos\theta)
- (L+r) P_{L+r-1}(\cos\theta).
\]
Dividing by $L+r+1$ and writing this relation in matrix form yields
\begin{equation}\label{MATRIZ_LEGEN1}
\begin{pmatrix}
P_{L+r+1} \\
P_{L+r}
\end{pmatrix}
=
\begin{pmatrix}
\dfrac{2L+2r+1}{L+r+1}\cos\theta &
-\dfrac{L+r}{L+r+1} \\
1 & 0
\end{pmatrix}
\begin{pmatrix}
P_{L+r} \\
P_{L+r-1}
\end{pmatrix}.
\end{equation}

By hypothesis, each polynomial $P_{L+k}$ can be expressed as
\[
P_{L+k}
= A_k^L P_L + B_k^L P_{L-1}.
\]
Substituting the expressions for $P_{L+r}$ and $P_{L+r-1}$ into
\eqref{MATRIZ_LEGEN1}, we obtain
\begin{equation}\label{EQmatrix1}
\begin{pmatrix}
P_{L+r+1} \\
P_{L+r}
\end{pmatrix}
=
\begin{pmatrix}
\dfrac{2L+2r+1}{L+r+1}\cos\theta &
-\dfrac{L+r}{L+r+1} \\
1 & 0
\end{pmatrix}
\begin{pmatrix}
A_r^L & B_r^L \\
A_{r-1}^L & B_{r-1}^L
\end{pmatrix}
\begin{pmatrix}
P_L \\
P_{L-1}
\end{pmatrix}.
\end{equation}

On the other hand, applying the representation of $P_{L+r+1}$ directly gives
\[
\begin{pmatrix}
P_{L+r+1} \\
P_{L+r}
\end{pmatrix}
=
\begin{pmatrix}
A_{r+1}^L & B_{r+1}^L \\
A_r^L & B_r^L
\end{pmatrix}
\begin{pmatrix}
P_L \\
P_{L-1}
\end{pmatrix}.
\]

Since the vectors $(P_L, P_{L-1})^{\mathsf T}$ are the same in both expressions,
equating the corresponding coefficient matrices in \eqref{EQmatrix1} yields
\[
\begin{pmatrix}
A_{r+1}^L & B_{r+1}^L \\
A_r^L & B_r^L
\end{pmatrix}
=
\begin{pmatrix}
\dfrac{2L+2r+1}{L+r+1}\cos\theta &
-\dfrac{L+r}{L+r+1} \\
1 & 0
\end{pmatrix}
\begin{pmatrix}
A_r^L & B_r^L \\
A_{r-1}^L & B_{r-1}^L
\end{pmatrix}.
\]
Comparing the entries of both sides gives the stated recurrence relations for
$A_{r+1}^L$ and $B_{r+1}^L$.
\end{proof}

\end{document}
